% Part: computability
% Chapter: computability-theory
% Section: universal-part-function

\documentclass[../../../include/open-logic-section]{subfiles}

\begin{document}

\olfileid{cmp}{thy}{uni}
\olsection{通用部分可计算函数}

\begin{thm}
\ollabel{thm:univ-comp}
存在通用部分可计算函数~$\fn{Un}(e,x)$。换言之，存在函数 $\fn{Un}(e,x)$ 满足：
\begin{enumerate}
\item $\fn{Un}(e,x)$ 是部分可计算的。
\item 若 $f(x)$ 是任意部分可计算函数，则存在自然数 $e$，使得对每个~$x$ 都有 $f(x) \simeq \fn{Un}(e,x)$。
\end{enumerate}
\end{thm}

\begin{proof}
令 $\fn{Un}(e,x) \simeq U(\umin{s}{T(e,x,s)})$，其中 $U$ 和 $T$ 如 Kleene的范式定理（\olref[nfm]{thm:normal-form}）所述。
\end{proof}

\begin{explain}
这只是对“部分可计算函数存在有效枚举”的精确表述；基本想法是，若用 $f_e(x) = \fn{Un}(e,x)$ 定义函数 $f_e$，则序列 $f_0$, $f_1$, $f_2$, \dots 包含了所有部分可计算函数，并且 $f_e(x)$ 关于 $e$ 和~$x$ 是“一致地”可计算的。为简单起见，我们使用了一个通用于一元函数的二元函数，但通过对数列编码，很容易将其推广到更多自变元的情形。例如，注意到若 $f(x,y,z)$ 是三元部分递归函数，则函数 $g(x) \simeq f((x)_0, (x)_1, (x)_2)$ 是一元递归函数。
\end{explain}

\end{document}